\section*{Глава 1. Анализ предметной области}
\addcontentsline{toc}{section}{Глава 1. Анализ предметной области}
\stepcounter{section}

\subsection{Обзор существующих программных продуктов}

Для определения места разрабатываемого приложения был проведён анализ существующих сервисов, 
работающих с цитатами из фильмов.

\textbf{PlayPhrase.me} — поисковик фраз из фильмов с возможностью воспроизведения видеоклипа.   			  \textit{Плюсы:} быстрый поиск, большая база. 					\textit{Минусы:} нет личных коллекций, нет аналитики \cite{playphrase}. 	 

\textbf{QuoDB / Subzin} — сервисы поиска по субтитрам с указанием таймкода.                                    \textit{Плюсы:} точность поиска.                             \textit{Минусы:} только поиск, без персонализации \cite{quodb}.  

\textbf{Film Buff} — мобильное приложение для ведения дневника фильмов, позволяет добавлять заметки и цитаты.  \textit{Плюсы:} личный кабинет. 								\textit{Минусы:} нет веб-версии, аналитика по цитатам отсутствует.

\textbf{Wikiquote} — вики-проект, содержащий коллекции цитат из фильмов, книг и других источников. 		       \textit{Плюсы:} структурированные страницы по каждому фильму. \textit{Минусы:} нет личного кабинета, нельзя сохранять свои цитаты \cite{wikiquote}. 	

\textbf{IMDb} — крупнейшая база данных о фильмах, содержит раздел «Quotes» для каждого фильма. 				  \textit{Плюсы:} огромная база. 								\textit{Минусы:} только просмотр, без возможности вести личную коллекцию \cite{imdb}. 

Результаты сравнения сведены в таблицу~\ref{tab:analogs}.

\begin{table}[ht!]
    \centering
    \caption{Сравнение существующих решений}
    \label{tab:analogs}
    \begin{tabular}{|p{0.25\textwidth}|p{0.2\textwidth}|p{0.2\textwidth}|p{0.2\textwidth}|}
        \hline
        \textbf{Сервис} 					& \textbf{Личная коллекция} & \textbf{Аналитика} & \textbf{Импорт/экспорт} \\ \hline
        PlayPhrase.me 						& Нет 						& Нет 				 & Нет 					   \\ \hline
        QuoDB 								& Нет 						& Нет 				 & Нет 					   \\ \hline
        Film Buff 							& Да (только мобильное) 	& Нет 				 & Нет 					   \\ \hline
        Wikiquote 							& Нет 						& Нет 				 & Нет 					   \\ \hline
        IMDb 								& Нет 						& Нет 				 & Нет 					   \\ \hline
        \textbf{Разрабатываемое приложение} & \textbf{Да} 				& \textbf{Да} 		 & \textbf{Да} 			   \\ \hline
    \end{tabular}
\end{table}

\textbf{Вывод:} существующие сервисы не предоставляют пользователю возможности вести личный 
цитатник с аналитикой и импортом/экспортом. Разрабатываемое приложение заполняет эту нишу.

\subsection{Анализ программных инструментов разработки}

Для реализации веб-приложения был проведён сравнительный анализ современных инструментов 
разработки по следующим критериям: простота освоения, производительность, наличие необходимых 
расширений, совместимость с выбранной архитектурой и достаточность функционала для решения 
поставленных задач.

\subsubsection{Выбор backend-фреймворка}

Были рассмотрены три наиболее популярных подхода на языке Python: Django, FastAPI и Flask. 
Сравнение представлено в таблице~\ref{tab:backend}.

\begin{table}[ht!]
    \centering
    \caption{Сравнение backend-фреймворков}
    \label{tab:backend}
    \begin{tabular}{|p{0.2\textwidth}|p{0.2\textwidth}|p{0.2\textwidth}|p{0.25\textwidth}|}
        \hline
        \textbf{Критерий}             & \textbf{Django}        & \textbf{FastAPI}               & \textbf{Flask}                \\ \hline
        Полнота функционала           & Высокая (всё включено) & Средняя (требует доп. модулей) & Низкая (требует доп. модулей) \\ \hline
        Порог вхождения               & Высокий                & Средний                        & Низкий                        \\ \hline
        Скорость                      & Средняя                & Высокая (асинхронность)        & Средняя                       \\ \hline
        Гибкость                      & Низкая                 & Высокая                        & Высокая                       \\ \hline
        Подходит для учебного проекта & Да, но избыточен       & Да                             & Да                            \\ \hline
    \end{tabular}
\end{table}

\textbf{Обоснование выбора Flask.} Django, несмотря на свою мощь, является тяжёлым решением 
для небольшого веб-приложения: он навязывает собственную структуру проекта и ORM, что 
ограничивает гибкость. FastAPI, хотя и обладает высокой производительностью за счёт 
асинхронности, требует понимания концепций async/await, что усложняет разработку для 
учебного проекта. Flask был выбран как минималистичный, но расширяемый фреймворк: он 
предоставляет только самое необходимое (маршрутизацию, обработку запросов), а все 
дополнительные функции (работа с БД, аутентификация, формы) подключаются через расширения. 
Это позволяет контролировать архитектуру и понимать каждый компонент \cite{flask_docs}. 

\subsubsection{Выбор базы данных}

Рассматривались две основные СУБД: SQLite и PostgreSQL. SQLite — встраиваемая файловая БД,
 удобная для прототипирования, но имеющая ограничения при работе с несколькими пользователями 
 одновременно и при выполнении сложных агрегирующих запросов. PostgreSQL — полноценная 
 серверная СУБД, поддерживающая расширенные типы данных, сложные запросы с GROUP BY, оконные 
 функции и транзакции, что критично для модуля аналитики \cite{postgresql_docs}. 

\textbf{Обоснование выбора PostgreSQL.} В разрабатываемом приложении планируется агрегация 
данных для построения графиков (подсчёт цитат по месяцам, топ фильмов, распределение по 
жанрам). PostgreSQL эффективно обрабатывает такие запросы благодаря оптимизатору и поддержке 
сложных агрегаций. Для взаимодействия с БД используется SQLAlchemy — мощный ORM, который 
позволяет работать с БД на уровне объектов Python, не теряя возможности писать сырой SQL 
при необходимости \cite{sqlalchemy_docs}. 

\subsubsection{Выбор инструмента для визуализации}

Для отображения статистических данных были рассмотрены две альтернативы: генерация графиков 
на сервере (Matplotlib) и построение графиков на стороне клиента (Chart.js).

Matplotlib позволяет строить графики на Python, сохраняя их в виде изображений. Однако этот 
подход имеет недостатки: графики получаются статичными, неинтерактивными, а их внешний вид 
требует дополнительной настройки. Кроме того, генерация изображений на сервере создаёт 
дополнительную нагрузку.

Chart.js — легковесная JavaScript-библиотека, которая рисует графики в браузере на основе 
переданных данных. Данные для графиков сервер отдаёт в формате JSON (результаты агрегирующих 
SQL-запросов), а Chart.js превращает их в интерактивные диаграммы (при наведении 
отображаются точные значения, графики анимированы).

\textbf{Обоснование выбора Chart.js.} Этот подход разгружает сервер, даёт современный 
интерактивный интерфейс и не требует сложной настройки. При этом агрегация данных 
по-прежнему выполняется на сервере с помощью SQL (или SQLAlchemy), что позволяет 
использовать всю мощь реляционной БД \cite{chartjs_docs}. 

\subsubsection{Выбор инструментов для аутентификации и фронтенда}

\textbf{Flask-Login.} Стандартное расширение для Flask, управляющее сессиями пользователей. 
Оно легко настраивается, интегрируется с любой моделью пользователя и предоставляет готовые 
декораторы (\texttt{@login\_required}) для защиты маршрутов \cite{flask_login_docs}. 

\textbf{Bootstrap 5.} Фреймворк для адаптивной вёрстки. Выбран из-за широкой распространённости, 
хорошей документации и наличия готовых компонентов (карточки, формы, модальные окна), 
что ускоряет разработку интерфейса \cite{bootstrap_docs}. 

\textbf{Работа с файлами.} Для импорта и экспорта цитат используются встроенные модули 
Python: \texttt{csv} для работы с CSV-файлами и \texttt{json} для работы с JSON. Это 
стандартные библиотеки, не требующие установки дополнительных пакетов.

\subsubsection{Итоговая таблица выбранного стека}

\begin{table}[ht!]
    \centering
    \caption{Итоговый стек технологий}
    \label{tab:stack}
    \begin{tabular}{|p{0.3\textwidth}|p{0.6\textwidth}|}
        \hline
        \textbf{Компонент} & \textbf{Выбранный инструмент}        \\ \hline
        Backend-фреймворк  & Flask                                \\ \hline
        База данных        & PostgreSQL + SQLAlchemy              \\ \hline
        Аутентификация     & Flask-Login                          \\ \hline
        Фронтенд           & Bootstrap 5                          \\ \hline
        Визуализация       & Chart.js                             \\ \hline
        Работа с файлами   & csv, json (встроенные модули Python) \\ \hline
    \end{tabular}
\end{table}

Таким образом, выбранный стек технологий оптимален для реализации всех функциональных 
требований курсового проекта: он обеспечивает необходимую производительность, гибкость и 
простоту разработки при сохранении высокого качества конечного продукта.

\subsection{Формулировка цели и задач работы}

На основе проведённого анализа сформулированы цель и задачи курсового проектирования.

\textbf{Цель работы} — разработать веб-приложение для ведения личного цитатника по фильмам, 
обеспечивающее пользователю возможность добавлять цитаты, управлять ими и получать наглядную 
статистику.

Для достижения цели необходимо решить следующие \textbf{задачи}:

\begin{enumerate}
    \item Провести анализ существующих аналогов и обосновать выбор технологического стека.
    \item Спроектировать базу данных для хранения информации о пользователях, фильмах, 
    жанрах и цитатах.
    \item Реализовать аутентификацию и авторизацию с разделением ролей (пользователь, 
    администратор).
    \item Разработать CRUD-интерфейс для управления фильмами и цитатами.
    \item Реализовать модуль аналитики: топ фильмов по количеству цитат, распределение 
    цитат по жанрам и другие статистические показатели.
    \item Добавить функции импорта цитат из файлов CSV/JSON и экспорта коллекции в эти 
    форматы.
    \item Выполнить тестирование разработанного приложения и подготовить документацию.
\end{enumerate}